\documentclass[12pt]{article}
\usepackage[T1]{fontenc}
\usepackage[utf8]{inputenc}
\usepackage{lmodern}
\usepackage{textcomp}
\usepackage{enumitem}
\usepackage{geometry}
\geometry{margin=1in}
\begin{document}
\begin{center}
Answer Key - Psychology - Undergraduate Level Assessment 8 \
\small (Answers are ordered exactly as in the question paper.)
\end{center}

\section*{Multiple Choice}
\begin{enumerate}[label=\arabic*.]
\item \textbf{Answer:} A
\\ \textit{Options:} A) blocked potential.; B) severe trauma.; C) genetic predisposition.; D) cognitive distortions.; E) lack of emotional intelligence.
\\ \textit{Note:} Practitioners of humanistic psychotherapy view psychopathology as stemming from the obstruction of an individual's inherent potential for growth and self-actualization, leading to psychological distress.
\item \textbf{Answer:} A
\\ \textit{Options:} A) receptionist; B) word processor; C) high school teacher; D) clinical psychologist.; E) software developer
\\ \textit{Note:} Vestibule training is appropriate for a receptionist role because it allows individuals to practice skills in a simulated environment that closely resembles their actual work setting, facilitating the development of necessary interpersonal and technical skills without the immediate pressures of a live workplace.
\item \textbf{Answer:} A
\\ \textit{Options:} A) Does not support self-efficacy for change; B) Empathy through reflective listening; C) Roll with resistance; D) Avoid direct confrontation
\\ \textit{Note:} The correct answer, "does not support self-efficacy for change," is accurate because Motivational Interviewing fundamentally aims to enhance a patient's belief in their ability to change, rather than undermine it.
\item \textbf{Answer:} D
\\ \textit{Options:} A) more effective for minority clients; B) only effective for nonminority clients; C) more effective when client and therapist have different racial/ethnic origins; D) equally effective; E) only effective when client and therapist speak the same language
\\ \textit{Note:} Research has demonstrated that psychotherapy yields similar positive outcomes for both minority and nonminority clients, highlighting the effectiveness of therapeutic techniques across diverse cultural backgrounds.
\item \textbf{Answer:} A
\\ \textit{Options:} A) itch; B) balance; C) heat; D) touch; E) pain
\\ \textit{Note:} Itch is considered a specific type of sensation that results from irritation of the skin rather than a basic somatosensation, which typically includes touch, pain, temperature, and proprioception.
\end{enumerate}

\section*{True / False}
\begin{enumerate}[label=\arabic*.]
\setcounter{enumi}{5}
\item \textbf{Answer:} True
\\ \textit{Options:} A) True; B) False
\\ \textit{Note:} Pascale's research into the processing strategies that children use to learn new information aligns with the focus of cognitive psychology, which examines mental processes such as perception, memory, and learning.
\item \textbf{Answer:} False
\\ \textit{Options:} A) True; B) False
\\ \textit{Note:} Children typically develop the fine motor skills required for clapping hands before they acquire the more complex dexterity needed to tie their shoes.
\item \textbf{Answer:} True
\\ \textit{Options:} A) True; B) False
\\ \textit{Note:} Neurotransmitters are indeed deactivated after they transmit their signals across the synapse through processes such as re-uptake, where they are taken back into the presynaptic neuron, or enzymatic breakdown, ensuring that their effects are temporary and allowing for precise regulation of neuronal communication.
\item \textbf{Answer:} False
\\ \textit{Options:} A) True; B) False
\\ \textit{Note:} Children's perception of letters during development is primarily influenced by visual and auditory stimuli, especially the sounds associated with letters, rather than sensory experiences related to taste and smell.
\item \textbf{Answer:} True
\\ \textit{Options:} A) True; B) False
\\ \textit{Note:} At birth, the visual system is least developed compared to other senses because the neural pathways and connections required for processing visual information are still maturing, leading to limited visual acuity and perception in newborns.
\end{enumerate}

\section*{Long Form Response}
\begin{enumerate}[label=\arabic*.]
\setcounter{enumi}{10}
\item \textbf{Answer:} In psychoanalytic theory, "catharsis" refers to the process of releasing and thereby providing relief from strong or repressed emotions, often achieved through therapeutic dialogue and exploration of unconscious conflicts. This emotional release is believed to alleviate psychological distress and anxiety by allowing individuals to confront and process their feelings. In contrast, behaviorists interpret catharsis through the lens of anxiety redirection, positing that the behaviors associated with cathartic experiences can serve as a means to redirect anxiety into more manageable expressions. Rather than focusing on emotional release, behaviorists would emphasize how therapeutic influences can modify behaviors and responses, thus helping individuals learn healthier coping mechanisms for anxiety. Ultimately, while psychoanalysts value the emotional catharsis itself, behaviorists highlight the importance of the behavioral changes that accompany this emotional process.
\\ \textit{Note:} correct because it accurately defines "catharsis" in psychoanalytic theory as the release of repressed emotions to alleviate anxiety and contrasts this with the behaviorist perspective, which focuses on redirecting anxiety through observable behaviors rather than emotional processing.
\item \textbf{Answer:} Double-bind messages occur when a communicator presents conflicting messages that create confusion for the recipient. In the scenario where a brother offers his car to his sister while referencing a past accident, he simultaneously conveys a willingness to help and an underlying caution. This creates a double-bind situation; his sister may feel encouraged to borrow the car but also anxious about the implications of her previous accident. The psychological impact of such conflicting signals can lead to feelings of guilt or inadequacy in the sister, as she may interpret the brother's concern as a lack of trust in her driving abilities. Ultimately, this type of communication can strain relationships and complicate decision-making processes.
\\ \textit{Note:} correct because it identifies the conflicting messages inherent in the brother's offer and the reference to the past accident, illustrating how this creates confusion and anxiety for the sister, thus exemplifying the concept of double-bind messages in communication.
\end{enumerate}

\vfill
\noindent\textit{End of Answer Key}
\end{document}